% Revised Methods and Results.
% Structural changes are summarized at the end of this file.
% Bracketed bold placeholders mark numbers the draft never reported; they must be
% filled from the logs before submission. Requires the booktabs package for Table I.

\section{Methods}

\subsection{Wall-Following Control}

The crawler holds a commanded standoff $d_{\mathrm{ref}}$ from the pipe wall while
traversing the pipe. Let the standoff error be
%
\begin{equation}
e(t) = \hat{d}(t) - d_{\mathrm{ref}},
\label{eq:error}
\end{equation}
%
where $\hat{d}(t)$ is the estimated distance from the crawler body to the pipe wall and
$d_{\mathrm{ref}}$ is the commanded standoff. The lateral correction that drives this
error to zero is given by
%
\begin{equation}
u(t) = \mathrm{sat}\!\left(k_p\, e(t) + k_d\, \dot{e}(t),\; u_{\max}\right),
\label{eq:pd}
\end{equation}
%
where $k_p$ and $k_d$ are the proportional and derivative gains,
$\mathrm{sat}(x, u_{\max}) = \min(\max(x, -u_{\max}), u_{\max})$ saturates the
correction, and $u_{\max}$ bounds the lateral speed the tracks may command. The
saturation keeps a large transient error from commanding a lateral rate that would
break track traction against the pipe wall. The correction enters the two track
velocities differentially, which are streamed to the drive controller at
$f_{\mathrm{ctrl}} = 100$~Hz through a real-time motor-control interface. Parameter
values are listed in Table~\ref{tab:params}.

The distance estimate $\hat{d}$ comes from a three-state extended Kalman filter driven
by two time-of-flight range measurements of the pipe wall, sampled at
$f_{\mathrm{ctrl}}$. The raw ranges are low-pass filtered before the update to reject
sensor noise, which at this standoff is comparable to the tolerance the controller must
hold. Process and measurement noise covariances $\mathbf{Q}$ and $\mathbf{R}$ are given
in Table~\ref{tab:params}.
% AUTHOR: the draft describes both a single-pole IIR pre-filter and an EKF without
% stating which one supplies the distance the PD law reads. Confirm against the source
% and, if the EKF output is unused by control, delete it from the paper entirely.
% AUTHOR: the three-state EKF and two measurements are inferred from the dimensions of
% Q and R in the draft. State the state vector explicitly if that inference is wrong.

We evaluate two controllers that differ only in how the range measurement is corrected
for sensor temperature. The \emph{compensated} controller corrects each range against a
reference sensor mounted at the front bulkhead. The \emph{uncompensated} controller uses
the raw range. Both share the estimator and the control law of
(\ref{eq:error})--(\ref{eq:pd}), so any difference between them is attributable to the
compensation alone.
% AUTHOR: the draft never states how the front-bulkhead reference sensor enters the
% range correction. One sentence of physical mechanism is needed here, otherwise the
% Results claim about compensation has no described cause.

\subsection{Experimental Setup}

Runs were conducted in a 12~m test pipe. At the start of each run the crawler was placed
at the pipe entrance and traversed the full length autonomously at 0.1~m/s. Wall
distance was logged at $f_{\mathrm{ctrl}}$ for the entire traverse, giving the
per-sample error of (\ref{eq:error}) used in all reported statistics.

We performed 16 runs, alternating between the compensated and uncompensated controller
so that any drift in pipe or ambient conditions over the campaign acted on both
controllers equally. The alternation yields eight paired runs.

\begin{table}[t]
\caption{Control and Estimation Parameters}
\label{tab:params}
\centering
\small
\setlength{\tabcolsep}{4pt}
\begin{tabular}{llr}
\toprule
Symbol & Description & Value \\
\midrule
$d_{\mathrm{ref}}$ & Commanded wall standoff (mm) & \textbf{[value]} \\
$k_p$ & Proportional gain, lateral correction (s$^{-1}$) & 2.1 \\
$k_d$ & Derivative gain, lateral correction (dimensionless) & 0.35 \\
$u_{\max}$ & Saturation limit on lateral correction (m/s) & 0.15 \\
$f_{\mathrm{ctrl}}$ & Sensing and command rate (Hz) & 100 \\
$\mathbf{Q}$ & Estimator process-noise covariance & $\mathrm{diag}(10^{-2}, 10^{-2}, 10^{-3})$ \\
$\mathbf{R}$ & Estimator measurement-noise covariance & $\mathrm{diag}(4.0, 4.0)$ \\
\bottomrule
\end{tabular}
\end{table}

\subsection{Metrics and Statistical Analysis}

The accuracy metric is the mean absolute standoff error $\overline{|e|}$ over a run.
Across runs, each controller is summarized as mean $\pm$ standard deviation of
$\overline{|e|}$. The two controllers are compared by the paired difference in
$\overline{|e|}$ across the eight run pairs, reported as a magnitude with a 95\%
confidence interval and tested with a paired $t$-test. All runs enter the statistics.
% AUTHOR: the draft reported 3.1 mm after excluding two runs that stayed outside a 5 mm
% band. Excluding the worst runs biases the headline metric, so all 16 runs must be
% included and the aggregate recomputed. The paired difference, its confidence interval,
% and the p-value were not in the draft and must be computed from the logs.

\section{Results}

The compensated controller held a mean absolute standoff error of
\textbf{[x.x]}~$\pm$~\textbf{[x.x]}~mm across the 12~m traverse, against
\textbf{[x.x]}~$\pm$~\textbf{[x.x]}~mm for the uncompensated controller
(Fig.~\ref{fig:error}). The paired difference is \textbf{[x.x]}~mm, 95\% CI
[\textbf{[x.x]}, \textbf{[x.x]}], a reduction of \textbf{[xx]}\%.
% AUTHOR: fill from the recomputed all-runs aggregate. The draft's 3.1 mm excluded two
% runs and reported no dispersion; the uncompensated aggregate was never reported at all,
% so the paper currently states a comparison it does not quantify.

The two controllers separate over the course of a traverse rather than uniformly
(Fig.~\ref{fig:error}). Within the first 30~s the uncompensated controller held the
lower error, by \textbf{[x.x]}~mm on average. Beyond 30~s its error grew for the
remainder of the traverse, reaching \textbf{[x.x]}~mm at the pipe exit, while the
compensated controller stayed within \textbf{[x.x]}~mm of its initial level.

\begin{figure}[t]
% AUTHOR: verify the line styles named below against the plotted figure.
\caption{Absolute standoff error along the 12~m traverse for the compensated and
uncompensated controllers. Solid lines are the mean across runs and shaded bands one
standard deviation, over eight runs per controller. The dashed horizontal line marks the
5~mm standoff tolerance.}
\label{fig:error}
\end{figure}

% ---------------------------------------------------------------------------
% What changed, and why
%
% METHODS
% 1. The four-step "First/Second/Third/Fourth" cycle narration was a transcription of the
%    software loop. Replaced with the control law itself first and the estimator second,
%    as a supporting module that feeds it.
% 2. Code identifiers removed: wall_follow, ctrl_mode="wall_follow", and the vendor
%    set_track_velocity() call, which is now the capability by role ("streamed to the
%    drive controller at 100 Hz through a real-time motor-control interface" -- add a
%    citation to the interface if it has a literature presence).
% 3. Numeric constants triaged: gains, saturation limit, rate, Q and R moved to
%    Table I as symbol/description/value; the IIR coefficient alpha=0.2, the 40-iteration
%    cap, the every-5-cycles divergence check, and P0 = 100*I were deleted as internal
%    implementation constants that support no claim.
% 4. Deleted: the unused simulation mode (it produced none of the reported results); the
%    assert-retract sentence about impedance control; the signpost sentence pointing at
%    Section IV.
% 5. Added: the two controllers as named experimental variables, and the protocol,
%    metric definition, and statistical analysis moved here out of Results.
% 6. Added one physical reason each for the saturation and the estimator, per the rule
%    that supporting machinery gets a model name plus a reason, not a tuning dump.
%
% RESULTS
% 7. The whole protocol paragraph (placement, operator start, traverse speed, logging,
%    alternation) moved to Methods.
% 8. Interpretation removed from Results: "which shows that the thermal compensation is
%    effective", "confirms our design choice of placing the reference sensor at the front
%    bulkhead", and "this matches the thermal time constant of the sensor mount" all
%    belong in the Discussion, where they answer the Introduction's gap.
% 9. The clockwise/counterclockwise signed split was dropped. Direction is not the
%    finding here, and signed per-direction means obscure the aggregate accuracy claim.
% 10. Convergence bookkeeping ("14 of 16 runs converged ... the two non-converging runs
%     are excluded") removed. The exclusion is also a validity problem, not a style one:
%     see the AUTHOR note in Methods.
% 11. "was better" replaced by a quantified difference; every comparison now carries a
%     magnitude.
% 12. The corroded-pipe sentence was removed from Results. It is a limitation, and it
%     stated what the lab lacked. Reframed for the Discussion:
%     "Evaluation was confined to clean pipe. Extending it to corroded pipe requires only
%     access to a decommissioned line, not a change to the controller."
% 13. Figure caption rewritten to be self-contained, not to open with "The", and to
%     decode the mean lines, the SD bands, the run count, and the tolerance line.
%     The 5 mm tolerance line needs a stated source (specification or requirement);
%     a threshold line drawn without justification invites a reviewer question.
% 14. "sixteen runs" -> "16 runs" (IEEE style uses numerals for 10 and above), and the
%     contraction "it's" is gone with the sentence that carried it.
% ---------------------------------------------------------------------------
